\section{Bölüm sonu çözümlü problemler}

\begin{enumerate}
  \item \textbf{Test istatistiği ve karar.} Gözlenen fark 4,5, sıfır hipotezi değeri 0 ve standart hata 1,5'tir. Test istatistiğini bulun; çift yönlü $p=0{,}006$ ise sonucu yazın.

  \textbf{Çözüm:} $t=(4{,}5-0)/1{,}5=3$ olur. $p=0{,}006<0{,}05$ olduğundan sıfır farkına karşı kanıt vardır. Sonuç $H_0$'ın doğru olma olasılığının yüzde 0,6 olduğu anlamına gelmez.

  \item \textbf{Hata türleri.} Gerçekte etkisiz bir programı etkili ilan etmek ve gerçekte etkili bir programı etkisiz saymak hangi hata türleridir?

  \textbf{Çözüm:} Etkisiz programı etkili ilan etmek yanlış pozitif, yani Tip I hatadır. Etkili program için yeterli kanıt bulamamak yanlış negatif, yani Tip II hatadır. Tip II hata olasılığı $\beta$, güç ise $1-\beta$'dır.

  \item \textbf{Çoklu test.} Bağımsız olduğu varsayılan 10 testin her biri $\alpha=0{,}05$ ile yürütülürse en az bir yanlış pozitif bulma olasılığını yaklaşık hesaplayın.

  \textbf{Çözüm:} $1-(1-0{,}05)^{10}=1-0{,}95^{10}\approx0{,}401$'dir. Olasılık yaklaşık yüzde 40,1'e yükselir. Testler bağımlı olduğunda kesin değer değişir; çoklu karşılaştırma sorunu yine sürer.
\end{enumerate}